\subsection{插值}

我们在区间 $[-1, 1]$ 上寻找由下式定义的函数 $f$ 的插值多项式：
\[
  f(x)=\frac{1}{1+8x^2}
\]

插值多项式为在 21 个等距点上构造的 \tkzimp{拉格朗日}多项式；次数为 $20$。

使用 \tkzname{tkz-fct} 时，多项式系数必须在别处计算，并将多项式写成 \tkzimp{Horner} 形式。使用 Lua 时，多项式直接由点计算：用 |tkzfn.fct| 将 Lua 函数存入表 |FN|。$f$ 的图像为红色，插值多项式为蓝色，使用的点为黄色。

\subsubsection{代码}
\begin{tkzexample}[code only]
\tkzFnDef{f}{1/(1+8*x^2)}
\directlua{
  local xs, ys = {}, {}
  for i = 0, 20 do
    xs[i] = -1 + i/10
    ys[i] = FN.f(xs[i])
  end
  FN.P = tkzfn.fct(function(x)
    local s = 0
    for i = 0, 20 do
      local l = 1
      for j = 0, 20 do
        if j == i then else l = l*(x-xs[j])/(xs[i]-xs[j]) end
      end
      s = s + ys[i]*l
    end
    return s
  end)
}
\begin{tikzpicture}
  \tkzFnInit[xmin=-1,xmax=1,ymin=-1.8,ymax=1.2,xstep=0.1,ystep=0.2]
  \tkzFnGrid
  \tkzFnAxes
  \tkzFnDraw[samples=400,line width=4pt,color=red,opacity=.5](f)
  \tkzFnDraw[samples=400,line width=1pt,color=blue](P)
  \tikzset{tkzfn point/.style={draw=black,fill=yellow}}
  \tkzFnMarkPoints[size=3pt](f){-1,-0.9,...,1}
\end{tikzpicture}
\end{tkzexample}

结果见下页，可以看到 \tkzimp{龙格}现象。

\subsubsection{图形}

\tkzFnDef{f}{1/(1+8*x^2)}
\directlua{
  local xs, ys = {}, {}
  for i = 0, 20 do
    xs[i] = -1 + i/10
    ys[i] = FN.f(xs[i])
  end
  FN.P = tkzfn.fct(function(x)
    local s = 0
    for i = 0, 20 do
      local l = 1
      for j = 0, 20 do
        if j == i then else l = l*(x-xs[j])/(xs[i]-xs[j]) end
      end
      s = s + ys[i]*l
    end
    return s
  end)
}
\begin{sidewaysfigure}[htbp]
\centering
\begin{tikzpicture}[scale=.75]
  \tkzFnInit[xmin=-1,xmax=1,ymin=-1.8,ymax=1.2,xstep=0.1,ystep=0.2]
  \tkzFnGrid
  \tkzFnAxes
  \tkzFnDraw[samples=400,line width=4pt,color=red,opacity=.5](f)
  \tkzFnDraw[samples=400,line width=1pt,color=blue](P)
  \tikzset{tkzfn point/.style={draw=black,fill=yellow}}
  \tkzFnMarkPoints[size=3pt](f){-1,-0.9,...,1}
\end{tikzpicture}
\caption{插值：$\dfrac{1}{1+8x^2}$}
\end{sidewaysfigure}

\endinput
